\documentclass{article}
\usepackage[utf8]{inputenc}
\usepackage[T1]{fontenc}
\usepackage{geometry}
\usepackage{graphicx}
\usepackage{float}
\usepackage{hyperref}
\usepackage[english,indonesian]{babel}

\title{Optimasi Pemrosesan Template LaTeX Menggunakan Pendekatan Berbasis AI}
\author{Academic Support Agent}
\date{\today}

\begin{document}
\maketitle

\begin{abstract}

\end{abstract}

\chapter{PENDAHULUAN}


\chapter{TINJAUAN PUSTAKA}


% Original content (commented out):
% \documentclass[a4paper,12pt]{article}
\usepackage[utf8]{inputenc}
\usepackage{graphicx}
\usepackage{amsmath}

\begin{document}

\begin{titlepage}

\begin{center}
    \begin{figure}[ht]
        \centering
        \includegraphics[height=2cm]{Friedrich-Alexander-Universität_Erlangen-Nürnberg_Logo_07.2022.eps}
    \end{figure}
\end{center}

\begin{center}
    \textbf{\huge Praktikum} \\[0.5cm]
    \textbf{\huge Titel 1}
    
    \vspace{2cm}
    
 %   \textbf{\Large } \\[0.5cm]
    \textbf{\Large Grundlegende Eigenschaften  } \\[0.5cm]
    \textbf{\Large von Polymeren} \\
    
    \vspace{3,5cm}
    
   \textbf{Gruppe 5:} \\[1cm]
    Name 1 \\[0.5cm]
    Name 2 \\[0.5cm]
    Name 3 \\[0.5cm]
    
   
    
\end{center}

\vspace{1cm}
  
 \begin{flushleft}
    \textbf{Praktikum: \underline{\hspace{1cm}day/month/year\hspace{1cm}}} \\[1cm]
    \textbf{1. Abgabe: \underline{\hspace{1cm}day/month/year\hspace{1cm}}} \\[0.5cm]
    \textbf{2. Abgabe: \underline{\hspace{4.5cm}}}
    \end{flushleft}

\end{titlepage}

\end{document}


\end{document}